% !TEX program = xelatex
\documentclass[12pt]{article}
\usepackage[margin=28mm]{geometry}
\usepackage{hyperref}
\usepackage{amsmath, amsthm, amssymb, amsfonts}
\usepackage{tikz-cd}
\usepackage{fontspec}
\usepackage{unicode-math}
\usepackage{xeCJK}
\setmainfont{Latin Modern Roman}
\setmonofont{Latin Modern Mono}
\setCJKmainfont{Noto Serif CJK JP}
\usepackage{listings}
\lstset{basicstyle=\ttfamily\footnotesize, columns=fullflexible, breaklines=true, showstringspaces=false}
\newtheorem{theorem}{定理}
\newtheorem{lemma}{補題}
\theoremstyle{definition}
\newtheorem{definition}{定義}
\newtheorem{remark}{注意}

\title{lean\_linear\_algebra\_task1 — 日本語ノート}
\author{TeX generated by ChatGPT-5 Pro\\Lean source generated by CODEX (OpenAI)}
\date{2025-09-10}

\begin{document}
\maketitle

\begin{abstract}
Bourbaki の精神（順序対と射影）: 構造は順序対 (x, y) と射影 π₁, π₂ によって
  特色付けられる。線形代数の文脈では直積 `V × W` は各成分への射影を通じて
  議論でき、等式は射影により「成分ごと」に確認できる。
\end{abstract}

\section*{生成元の明示（Provenance)}
\begin{itemize}
  \item \textbf{TeX 生成}: ChatGPT-5 Pro
  \item \textbf{Lean ソース生成}: CODEX (OpenAI)
\end{itemize}

\section*{概説}
本ノートはアップロードされた Lean ソースをもとに、著者名に「TeXはChatGPT-5 Proで生成」「LeanファイルはCODEX (OpenAI)で生成」を明示し、
基本的な可換図式（用途に応じて積・線形代数の図式を含む）と、定理・定義の先頭行一覧、ソース全文を掲載します。

\section*{サマリー}
Total lines: 60\\
Defs: 0\\
Lemmas: 0\\
Theorems: 1\\
Structures: 0\\
Inductives: 0

\section*{図式（例）}

\[
\begin{tikzcd}
A \arrow[r, "f"] \arrow[d, "g"'] & B \arrow[d, "h"] \\
C \arrow[r, "k"'] & D
\end{tikzcd}
\]



\[
\begin{tikzcd}
V \arrow[r, "A"] \arrow[dr, swap, "B\circ A"] & W \arrow[d, "B"] \\
& X
\end{tikzcd}
\qquad
\begin{tikzcd}
V \arrow[r, "A_1"] \arrow[d, "A_2"'] & W \\
U \arrow[ur, swap, "A_1\circ A_2"] &
\end{tikzcd}
\]


\section*{定理・定義の一覧（先頭行）}
\begin{center}
\begin{tabular}{lll}
\textbf{種別} & \textbf{名前} & \textbf{シグネチャ} \\ \hline
theorem & \texttt{scalar\_mul\_distribute} & \texttt{{R V : Type*} [Semiring R] [AddCommMonoid V] [Module R V] (a : R) (v w : V) : a • (v + w) = a • v + a • w := by} \\
\end{tabular}
\end{center}

\section*{Lean ソース（全文）}
\begin{lstlisting}
import Mathlib.Algebra.Module.Basic
import Mathlib.Algebra.Module.Prod

/-!
  Bourbaki の精神（順序対と射影）: 構造は順序対 (x, y) と射影 π₁, π₂ によって
  特色付けられる。線形代数の文脈では直積 `V × W` は各成分への射影を通じて
  議論でき、等式は射影により「成分ごと」に確認できる。
-/

noncomputable section

namespace MyProjects.LinearAlgebra

-- ベクトル空間における分配法則（スカラー倍の加法に対する分配）
-- R: スカラーの型（半環）
-- V: R-加群（ベクトル空間を含む一般の加群）
theorem scalar_mul_distribute {R V : Type*}
    [Semiring R] [AddCommMonoid V] [Module R V]
    (a : R) (v w : V) :
    a • (v + w) = a • v + a • w := by
  simpa using (smul_add a v w)

/-!
  以降は「順序対と射影」による構造化の最小限の具体化。
  `V × W` 上での等式を、射影 `Prod.fst`（π₁）と `Prod.snd`（π₂）により
  成分ごとに確認する。
-/

section Projections

variable {R V W : Type*}
variable [Semiring R]
variable [AddCommMonoid V] [AddCommMonoid W]
variable [Module R V] [Module R W]

/-- π₁ による成分等式：`a • (x + y)` の第1成分は分配する。 -/
@[simp] lemma fst_smul_add (a : R) (x y : V × W) :
    (a • (x + y)).1 = a • x.1 + a • y.1 := by
  -- `V × W` での分配 `smul_add` を射影 `Prod.fst` に通すだけ
  simpa using congrArg Prod.fst (smul_add a x y)

/-- π₂ による成分等式：`a • (x + y)` の第2成分は分配する。 -/
@[simp] lemma snd_smul_add (a : R) (x y : V × W) :
    (a • (x + y)).2 = a • x.2 + a • y.2 := by
  simpa using congrArg Prod.snd (smul_add a x y)

-- クイックサニティチェック（型クラス前提のまま最小の `example`）
section Examples
  variable (a : R) (v w : V) (x y : V × W)
  example : a • (v + w) = a • v + a • w := by
    simpa using (smul_add a v w)
  example : (a • (x + y)).1 = a • x.1 + a • y.1 := by
    simpa using (fst_smul_add (R:=R) (V:=V) (W:=W) a x y)
  example : (a • (x + y)).2 = a • x.2 + a • y.2 := by
    simpa using (snd_smul_add (R:=R) (V:=V) (W:=W) a x y)
end Examples

end Projections

end MyProjects.LinearAlgebra

\end{lstlisting}

\end{document}
